\begin{frame}{Conclusiones}
\small
\begin{enumerate}
  \item Los estados financieros son la base del análisis.
  \item La liquidez debe revisar calidad de activos corrientes.
  \item La gestión mide eficiencia del ciclo operativo.
  \item La solvencia muestra dependencia de terceros.
  \item La rentabilidad exige comparar ventas, activos y patrimonio.
  \item El caso integrador conecta cálculo e interpretación.
  \item Las decisiones deben considerar contexto y metas.
  \item Ingeniería de Sistemas aporta automatización, trazabilidad y visualización.
\end{enumerate}
\end{frame}

\begin{frame}{Referencias APA 7}
\scriptsize
IFRS Foundation. (2018). \textit{Conceptual Framework for Financial Reporting}. https://www.ifrs.org/issued-standards/list-of-standards/conceptual-framework/

\vspace{0.4em}
Ministerio de Economía y Finanzas del Perú. (2019). \textit{Plan Contable General Empresarial modificado 2019}. https://www.gob.pe/92758-ministerio-de-economia-y-finanzas-pcge

\vspace{0.4em}
Superintendencia del Mercado de Valores. (2026). \textit{Información financiera}. https://www.gob.pe/institucion/smv/tema/informacion-financiera

\vspace{0.4em}
Amat, O. (2008). \textit{Análisis de estados financieros}. Gestión 2000.

\vspace{0.4em}
Warren, C. S., Reeve, J. M., \& Duchac, J. E. (2016). \textit{Contabilidad financiera}. Cengage Learning.
\end{frame}
